\chapter{The special-functions ladder}

\begin{orient}
\textbf{What this chapter is for.} Past a certain level of difficulty a definite integral simply does not have an elementary answer, and continuing to hunt for one is the mistake. The skill changes. Instead of searching for an antiderivative you ask which named function or named constant the integral is a disguised instance of, and then you work with that object's functional equations. Those equations are the real content: the reflection formula $\Gamma(s)\Gamma(1-s)=\pi/\sin\pi s$ evaluates an entire family of half-line integrals in one line, the recurrence $\psi(z+1)=\psi(z)+1/z$ turns a harmonic sum into a closed form, and the derivative relation $\Li_{s+1}'(x)=\Li_s(x)/x$ climbs a whole ladder of logarithmic integrals. This chapter builds that ladder rung by rung, starting from Gamma and Beta, passing through digamma and the polylogarithm, and ending at the elliptic integrals, where even the special functions of one variable give out. For each rung you get three things: the defining integral, the two or three identities that actually get used in practice, and the recognition cue that tells you which rung you are standing on. At the end you should be able to look at an unfamiliar definite integral and say, before computing anything, which named object its answer will be built from.
\end{orient}

\section{What a closed form is, past the elementary}

A student trained only on antiderivatives has a single test for success: can I write the answer with $\exp$, $\ln$, powers, and the trigonometric functions? By that test $\int_0^1\eu^{-x^{2}}\dd x$ has no answer, $\int_0^{\infty}\frac{\sin x}{x}\dd x$ has no answer, and $\int_0^1\frac{\dd x}{\sqrt{1-x^{4}}}$ has no answer. All three statements are false in any useful sense. The first equals $\tfrac{\sqrt\pi}{2}\erf(1)$, the second equals $\tfrac{\pi}{2}$, and the third equals $\Gamma(\tfrac14)^{2}/(4\sqrt{2\pi})$. What has happened is that the alphabet of acceptable answers grew.

The right way to think about the enlargement is this. A special function is a function defined by an integral or a series that is important enough to have been tabulated, named, and equipped with identities. Once it has identities, it can be computed with, and computing with it is exactly what integration at this level consists of. The functions worth carrying are few, and they come in a definite order of generality.

\begin{center}
\footnotesize
\setlength{\tabcolsep}{5pt}
\renewcommand{\arraystretch}{1.32}
\begin{tabular}{@{}llll@{}}
\toprule
\mh{Rung} & \mh{Object} & \mh{Defining integral} & \mh{The identity that does the work} \\
\midrule
1 & $\Gamma(s)$ & $\int_0^{\infty}x^{s-1}\eu^{-x}\dd x$ & $\Gamma(s+1)=s\Gamma(s)$ \\
2 & $\Gamma$ reflection & (same) & $\Gamma(s)\Gamma(1-s)=\pi/\sin\pi s$ \\
3 & $\mathrm{B}(p,q)$ & $\int_0^1t^{p-1}(1-t)^{q-1}\dd t$ & $\mathrm{B}=\Gamma\Gamma/\Gamma$ \\
4 & $\psi(z)$ & $\Gamma'/\Gamma$ & $\psi(z+1)=\psi(z)+\tfrac1z$ \\
5 & $\Li_s(x)$ & $\int_0^x\Li_{s-1}(t)\tfrac{\dd t}{t}$ & $\Li_s(1)=\zeta(s)$; Landen \\
6 & $\erf$, Fresnel & $\int_0^z\eu^{-t^{2}}\dd t$ & Complete the square \\
7 & $\Si$, $\Ci$ & $\int_0^z\tfrac{\sin t}{t}\dd t$ & $\Si(\infty)=\tfrac{\pi}{2}$ \\
8 & $\operatorname{Cl}_2(\theta)$ & $-\int_0^{\theta}\ln\abs{2\sin\tfrac t2}\dd t$ & Duplication in $\theta$ \\
9 & $K$, $E$ & $\int_0^{\pi/2}(1-k^{2}\sin^{2}\theta)^{\mp1/2}\dd\theta$ & Landen and AGM \\
\bottomrule
\end{tabular}
\end{center}

One unifying remark before the details. Almost every entry in that table is a \emph{Mellin transform}, that is, an integral of the form $\int_0^{\infty}x^{s-1}f(x)\dd x$ with $s$ a free parameter. Gamma is the Mellin transform of $\eu^{-x}$, Beta is the Mellin transform of $(1+x)^{-1}$ up to a change of variable, the zeta representation is the Mellin transform of $(\eu^{x}-1)^{-1}$, and the Fresnel values are the Mellin transform of $\sin$ at $s=\tfrac12$. That is why the same three moves keep working: substitute to expose a pure power against a fixed kernel, read off the parameter $s$, and check the strip of $s$ for which the integral converges. If you internalise nothing else from this chapter, internalise the habit of asking ``what is the free power, and what is the kernel''.

The ordering is not arbitrary. Rungs 1 through 4 are all the same function seen from different angles, since Beta is a ratio of Gammas and digamma is a logarithmic derivative of Gamma. Rung 5 is the systematic home of every integral built from $\ln$ against a rational function. Rungs 6 and 7 are what oscillation and Gaussian decay produce when the interval is cut short. Rung 8 is trigonometric logarithms. Rung 9 is the first genuinely new object, because $K$ and $E$ are functions of a modulus that no Gamma identity evaluates except at special points.

\begin{keynote}[Recognition, not derivation]
The practical skill here is almost entirely recognition. Nobody derives the reflection formula in the middle of a problem. What you do is look at $\int_0^{\infty}\frac{x^{-2/3}}{1+x}\dd x$, see a fractional power against a self-reciprocal denominator on the half-line, and write down $\pi/\sin(\pi/3)$. The identities go on a card; the cues go into your reflexes. When a problem resists, the productive question is not ``what substitution have I missed'' but ``which named constant is this answer going to contain''. If you can guess the constant, you can usually find the route backwards from it.
\end{keynote}

One more orientation point, about how these objects are actually used. Almost no real problem is a single lookup. A typical hard integral is a chain: a substitution turns it into a Beta function, the Beta function becomes a Gamma quotient, one of the Gammas has argument $\tfrac34$, reflection converts that into $\Gamma(\tfrac14)$, and the answer comes out as $\Gamma(\tfrac14)^{2}$ over something. Three named identities, none of them hard, applied in sequence. The difficulty in this material is never any single step; it is knowing which of the four or five available identities to apply next, and that is decided by looking at what is blocking you, not by looking at the original integral.

\section{Gamma, and the three identities that evaluate everything}

Gamma is where every route through this material begins, and it is worth being clear about why one function is doing so much work. The reason is that $\eu^{-x}$ against a power is the unique combination that is stable under both scaling and the recurrence: substituting $x\mapsto\lambda x$ changes the integral only by a power of $\lambda$, and integrating by parts changes the parameter only by one. Those two invariances are what turn a single integral into a function with a functional equation, and a functional equation is what lets you evaluate rather than merely define. Every other rung of the ladder inherits the same two properties in some form.

The Gamma function is the interpolation of the factorial, and for our purposes it is a machine that eats any integral of the shape ``power times exponential of a power on the half-line''. That is a much larger family than it sounds, because substitutions turn radicals, logarithms and Gaussians into that shape routinely. Start with the machine itself.

\tech{Gamma-function moment}{core}
\cue Any power of $x$ multiplied by an exponential of a power of $x$, integrated over $(0,\infty)$: $\int_0^{\infty}x^{m}\eu^{-\lambda x^{k}}\dd x$. The disguises are a base other than $\eu$, as in $2^{-x^{2}}$, and a radical, as in $\eu^{-\sqrt x}$, which is the case $k=\tfrac12$.
\method Substitute $u=\lambda x^{k}$, which converts everything into the defining integral $\Gamma(s)=\int_0^{\infty}u^{s-1}\eu^{-u}\dd u$. The result is worth memorising in one line:
\[\int_0^{\infty}x^{m}\eu^{-\lambda x^{k}}\dd x=\frac{1}{k\,\lambda^{(m+1)/k}}\,\Gamma\!\Bigl(\frac{m+1}{k}\Bigr),\qquad \frac{m+1}{k}>0,\ \lambda>0.\]
Then reduce the argument with $\Gamma(z+1)=z\Gamma(z)$ until it lands in $(0,1]$, and finish with $\Gamma(n)=(n-1)!$ or $\Gamma(\tfrac12)=\sqrt\pi$. A base other than $\eu$ is handled by $a^{-x^{k}}=\eu^{-x^{k}\ln a}$, so $\lambda=\ln a$.

\begin{worked}
Two instances, the second showing where the base goes.
\[\int_0^{\infty}x^{4}\eu^{-x^{3}}\dd x=\frac{1}{3}\Gamma\!\Bigl(\frac{5}{3}\Bigr)\approx0.300915,\]
directly from the formula with $m=4$, $k=3$, $\lambda=1$. Reduce further if you like: $\Gamma(\tfrac53)=\tfrac23\Gamma(\tfrac23)$.

Now $\displaystyle\int_{-\infty}^{\infty}\abs{x}^{3}2^{-x^{2}}\dd x$. The integrand is even, so this is $2\int_0^{\infty}x^{3}\eu^{-x^{2}\ln2}\dd x$. Here $m=3$, $k=2$, $\lambda=\ln2$, so $(m+1)/k=2$ and the formula gives
\[2\cdot\frac{1}{2(\ln2)^{2}}\Gamma(2)=\frac{1}{\ln^{2}2}\approx2.081369.\]
The exponent on $\ln2$ is $2$, not $1$, because the scale enters as $\lambda^{(m+1)/k}$.
\end{worked}

\begin{trap}
Two failures, both cheap to avoid. First, a non-$\eu$ base contributes $\ln a$ to \emph{every} power of the scale, so $(\ln2)^{2}$ downstairs in the example above, not $\ln2$; check the exponent by dimensional reasoning on $x$. Second, $\Gamma$ demands $\tfrac{m+1}{k}>0$. If it is not, the integral diverges at the origin and the formula still produces a finite-looking number, which is the worst possible failure mode. Test both endpoints before you use the answer: $x^{m}$ at $0$ and the exponential at $\infty$.
\end{trap}

\begin{seealsob}The Beta function; Gaussian and error function; Euler reflection and the Mellin transform; Zeta values and Apery's constant.\end{seealsob}

\begin{keynote}[Reducing a Gamma argument]
Whatever comes out of a substitution, the argument of $\Gamma$ can be brought to a standard place by three moves, applied in this order. Shift by integers with $\Gamma(z+1)=z\Gamma(z)$ until the argument lies in $(0,1]$. If it is now $\tfrac12$, you are done: $\Gamma(\tfrac12)=\sqrt\pi$. If it is a rational $\tfrac pq$ and $\tfrac{q-p}{q}$ also appears, pair them with reflection to get $\pi/\sin(\pi p/q)$. If both $z$ and $z+\tfrac12$ appear, use duplication. Anything left over, such as a lone $\Gamma(\tfrac13)$ or $\Gamma(\tfrac14)$, does not reduce further and is part of the answer. Recognising that last case early saves more time than any other habit in this chapter.
\end{keynote}

The recurrence $\Gamma(z+1)=z\Gamma(z)$ moves the argument by integers, which is enough to reduce any rational argument to a representative in $(0,1]$ but not enough to \emph{evaluate} that representative. Two further identities do the evaluating, and between them they close off the whole family of half-line integrals with fractional powers. The first is the reflection formula, which is by a wide margin the most valuable single fact in this chapter.

\tech{Euler reflection and the Mellin transform}{hard}
\cue A fractional power $x^{s-1}$ against a denominator like $1+x^{n}$ on $(0,\infty)$; equivalently, an integrand invariant in shape under $x\mapsto1/x$ with a non-integer exponent floating on top. Also any integral you have already reduced to $\Gamma(\alpha)\Gamma(1-\alpha)$.
\method The reflection formula is
\[\Gamma(s)\Gamma(1-s)=\frac{\pi}{\sin\pi s},\qquad 0<s<1,\]
and combining it with the Beta evaluation of the Mellin transform of $(1+x^{n})^{-1}$ gives the master formula
\[\int_0^{\infty}\frac{x^{s-1}}{1+x^{n}}\dd x=\frac{\pi}{n\sin(\pi s/n)},\qquad 0<s<n.\]
Everything on the half-line with a single self-reciprocal denominator is a case of this. For a compound denominator, partial-fraction it into pieces of this shape first, then apply the formula termwise.

\begin{worked}
The plain case first:
\[\int_0^{\infty}\frac{\dd x}{1+x^{3}}=\frac{\pi}{3\sin(\pi/3)}=\frac{\pi}{3}\cdot\frac{2}{\sqrt3}=\frac{2\pi}{3\sqrt3}\approx1.209200,\]
using $s=1$, $n=3$. A fractional numerator is no harder:
\[\int_0^{\infty}\frac{x^{-2/3}}{1+x}\dd x=\frac{\pi}{\sin(\pi/3)}=\frac{2\pi}{\sqrt3}\approx3.627599,\]
with $s=\tfrac13$, $n=1$. Notice that neither integral has an elementary antiderivative worth writing down; the first one does, by partial fractions, and it takes half a page.

The same formula run backwards evaluates a Gamma product. Since $\int_0^1(1-x^{3})^{-1/3}\dd x=\tfrac13\mathrm{B}(\tfrac13,\tfrac23)=\tfrac13\Gamma(\tfrac13)\Gamma(\tfrac23)$, and reflection gives $\Gamma(\tfrac13)\Gamma(\tfrac23)=\pi/\sin(\pi/3)$, the value is again $2\pi/(3\sqrt3)$.
\end{worked}

\begin{worked}[A denominator that is not $1+x^{n}$]
The same Beta-plus-reflection route handles any power of the denominator. Substituting $t=\tfrac{x}{1+x}$, or equivalently reading the Mellin transform of $(1+x)^{-\mu}$ directly,
\[\int_0^{\infty}\frac{x^{s-1}}{(1+x)^{\mu}}\dd x=\mathrm{B}(s,\mu-s)=\frac{\Gamma(s)\Gamma(\mu-s)}{\Gamma(\mu)},\qquad 0<s<\mu.\]
At $\mu=2$ this gives $\Gamma(s)\Gamma(2-s)=(1-s)\Gamma(s)\Gamma(1-s)=\pi(1-s)/\sin\pi s$, so
\[\int_0^{\infty}\frac{x^{-0.4}}{(1+x)^{2}}\dd x=\frac{0.4\pi}{\sin(0.6\pi)}\approx1.3213064,\]
confirmed by quadrature to eleven digits. The recurrence did the shifting and reflection did the evaluating, which is the standard division of labour.
\end{worked}

\begin{trap}
The convergence strip is the whole hypothesis and it is invisible in the formula. Outside $0<s<n$ the right-hand side still evaluates to a perfectly pleasant number while the integral diverges. Check the two endpoints separately: near $0$ the integrand behaves like $x^{s-1}$, which needs $s>0$; near $\infty$ it behaves like $x^{s-1-n}$, which needs $s<n$. A second trap: the formula is for $1+x^{n}$, not $1-x^{n}$. The latter has a pole on the path and its principal value is $\frac{\pi}{n}\cot(\pi s/n)$, a different function entirely.
\end{trap}

\begin{seealsob}The Beta function; Gamma-function moment; Legendre duplication; Quotable residue results.\end{seealsob}

Reflection relates $\Gamma$ at $s$ to $\Gamma$ at $1-s$. The other structural identity relates $\Gamma$ at $z$ and $z+\tfrac12$ to $\Gamma$ at $2z$, and it is what you reach for whenever half-integer arguments appear alongside integer ones. It is also the cleanest illustration in the chapter of a functional equation being \emph{produced} by a symmetry substitution rather than quoted.

\tech[Legendre duplication]{Legendre duplication, $\Gamma(z)\Gamma(z+\tfrac12)$}{medium}
\cue A Gamma quotient in which both $z$ and $z+\tfrac12$ appear, or a Beta function with equal arguments $\mathrm{B}(p,p)$, or a central binomial coefficient $\binom{2n}{n}$ that you want as a Gamma ratio. Also any Wallis-type product you want in closed form.
\method The identity is
\[\Gamma(z)\Gamma\!\Bigl(z+\frac12\Bigr)=2^{1-2z}\sqrt{\pi}\,\Gamma(2z).\]
Use it in whichever direction removes the half-integer. The general multiplication theorem
\[\prod_{j=0}^{n-1}\Gamma\Bigl(z+\frac jn\Bigr)=(2\pi)^{(n-1)/2}n^{\frac12-nz}\Gamma(nz)\]
is the same statement for $n$ factors, and the case $n=2$ is duplication.

\begin{worked}
Duplication is exactly the statement that $\mathrm{B}(p,p)$ has a symmetry. Compute
\[\mathrm{B}(p,p)=\int_0^1\bigl[x(1-x)\bigr]^{p-1}\dd x\]
by recentring at the midpoint: put $x=\tfrac{1+t}{2}$, so $x(1-x)=\tfrac{1-t^{2}}{4}$ and $\dd x=\tfrac{\dd t}{2}$, giving
\[\mathrm{B}(p,p)=\frac{4^{1-p}}{2}\int_{-1}^{1}(1-t^{2})^{p-1}\dd t=4^{1-p}\int_0^1(1-t^{2})^{p-1}\dd t=4^{1-p}\cdot\frac{1}{2}\mathrm{B}\Bigl(\frac12,p\Bigr),\]
the last step by $u=t^{2}$. Writing both Betas as Gamma quotients,
\[\frac{\Gamma(p)^{2}}{\Gamma(2p)}=\frac{2^{1-2p}\sqrt\pi\,\Gamma(p)}{\Gamma(p+\tfrac12)},\]
which rearranges to the duplication formula. Checked at $p=0.63$: both sides equal $1.3386477$.
\end{worked}

\begin{trap}
The factor $2^{1-2z}$, not $2^{-2z}$ and not $2^{2-2z}$. It is pinned by the single test $z=\tfrac12$, where the identity reads $\Gamma(\tfrac12)\Gamma(1)=2^{0}\sqrt\pi\,\Gamma(1)$, so evaluate at $z=\tfrac12$ every time you write it from memory rather than trusting the exponent.
\end{trap}

\begin{seealsob}The Beta function; Euler reflection and the Mellin transform; Reduction formulae and Wallis.\end{seealsob}

Reflection and duplication are identities about $\Gamma$ alone. The Beta function is what you actually meet, because the natural home of a two-parameter power weight is the unit interval and not the half-line. Beta is a ratio of Gammas and therefore adds no new information, but it adds an enormous amount of recognition value: once you can spot a Beta integrand, the Gamma machinery applies to bounded intervals too.

\tech{The Beta function}{core}
\cue A weight $x^{p-1}(1-x)^{q-1}$ on $[0,1]$, possibly after a substitution that maps your interval to $[0,1]$; or $\sin^{m}x\cos^{n}x$ on $[0,\tfrac{\pi}{2}]$ with exponents that are not both non-negative integers. A radical such as $\sqrt{1-x}$ or $(1-x^{3})^{-1/3}$ is a Beta weight with a fractional exponent.
\method Two forms, and the relation between them is the whole technique:
\[\mathrm{B}(p,q)=\int_0^1t^{p-1}(1-t)^{q-1}\dd t=\frac{\Gamma(p)\Gamma(q)}{\Gamma(p+q)},\qquad
\int_0^{\pi/2}\sin^{2p-1}x\cos^{2q-1}x\dd x=\frac12\mathrm{B}(p,q).\]
The trigonometric form follows from the algebraic one by $t=\sin^{2}x$. For a general power, $\int_0^1x^{a}(1-x^{b})^{c}\dd x=\tfrac1b\mathrm{B}\bigl(\tfrac{a+1}{b},c+1\bigr)$ via $u=x^{b}$; this one substitution covers most of what you will meet. Half-integer Gammas then come down to $\Gamma(\tfrac12)=\sqrt\pi$ by the recurrence.

\begin{worked}
An algebraic case and a trigonometric one.
\[\int_0^1x^{3}\sqrt{1-x}\dd x=\mathrm{B}\Bigl(4,\frac32\Bigr)=\frac{\Gamma(4)\Gamma(\tfrac32)}{\Gamma(\tfrac{11}{2})}
=\frac{6\cdot\tfrac{\sqrt\pi}{2}}{\tfrac{945}{32}\sqrt\pi}=\frac{32}{315}\approx0.101587,\]
using $\Gamma(\tfrac{11}{2})=\tfrac92\cdot\tfrac72\cdot\tfrac52\cdot\tfrac32\cdot\tfrac12\sqrt\pi=\tfrac{945}{32}\sqrt\pi$.

Now a fractional trigonometric exponent, where no elementary method applies at all:
\[\int_0^{\pi/2}\sqrt{\sin x}\dd x.\]
Match $2p-1=\tfrac12$ and $2q-1=0$, so $p=\tfrac34$, $q=\tfrac12$, and the value is
\[\frac12\mathrm{B}\Bigl(\frac34,\frac12\Bigr)=\frac{\Gamma(\tfrac34)\Gamma(\tfrac12)}{2\Gamma(\tfrac54)}
=\frac{\Gamma(\tfrac34)\sqrt\pi}{2\cdot\tfrac14\Gamma(\tfrac14)}=\frac{2\sqrt\pi\,\Gamma(\tfrac34)}{\Gamma(\tfrac14)}.\]
Reflection turns $\Gamma(\tfrac34)$ into $\pi\sqrt2/\Gamma(\tfrac14)$, so the answer is $2\sqrt{2\pi^{3}}/\Gamma(\tfrac14)^{2}\approx1.198140$. Two named identities, no antiderivative.
\end{worked}

\begin{worked}[The general power form, which covers most sightings]
One substitution handles the whole family $x^{m}(1+x^{n})^{-\mu}$ on the half-line. Put $u=\tfrac{x^{n}}{1+x^{n}}$, or simply combine the two previous results, to get
\[\int_0^{\infty}\frac{x^{m}}{(1+x^{n})^{\mu}}\dd x=\frac1n\mathrm{B}\Bigl(\frac{m+1}{n},\ \mu-\frac{m+1}{n}\Bigr).\]
With $m=2$, $n=3$, $\mu=\tfrac52$ this is $\tfrac13\mathrm{B}(1,\tfrac32)=\tfrac13\cdot\tfrac23=\tfrac29\approx0.222222$, matching quadrature. The case $\mu=1$ recovers the reflection master formula, so this one line contains everything in this section as a special case; memorise it in preference to the special cases.
\end{worked}

\begin{trap}
The exponent offsets, which are the single commonest error in this whole group. It is $p-1$ and $q-1$ in the algebraic form but $2p-1$ and $2q-1$ in the trigonometric form. Reading $\int_0^{\pi/2}\sin^{4}x\dd x$ as $\tfrac12\mathrm{B}(4,\tfrac12)$ instead of $\tfrac12\mathrm{B}(\tfrac52,\tfrac12)$ gives a clean wrong answer. Write down the two matching equations $2p-1=m$, $2q-1=n$ explicitly every time; do not do it in your head.
\end{trap}

\begin{seealsob}Gamma-function moment; Legendre duplication; Reduction formulae and Wallis; Elliptic integrals $K$ and $E$.\end{seealsob}

\section{Digamma: what happens when you differentiate in the parameter}

Gamma and Beta evaluate integrals whose integrands are pure powers. The moment a logarithm appears, differentiate the parameter: $\partial_p x^{p-1}=x^{p-1}\ln x$, so every $\ln x$ in the integrand is one $p$-derivative of a Beta or Gamma integral. Differentiating $\Gamma$ produces $\Gamma'$, and the useful normalisation is the logarithmic derivative, because that is what appears when you differentiate a Beta \emph{quotient}.

\tech{Digamma and polygamma}{hard}
\cue A logarithm multiplying a power weight; a harmonic number $H_n$ in a series you are trying to sum; the combination $\dfrac{1-x^{a-1}}{1-x}$ on $[0,1]$; or an answer that visibly ``wants'' to be $-\gamma$ plus something clean.
\method Define $\psi(z)=\Gamma'(z)/\Gamma(z)$. Four facts do all the work:
\[\psi(1)=-\gamma,\qquad \psi(z+1)=\psi(z)+\frac1z,\qquad \psi(1-z)-\psi(z)=\pi\cot\pi z,\]
together with the integral representation
\[\int_0^1\frac{1-x^{a-1}}{1-x}\dd x=\psi(a)+\gamma.\]
The recurrence gives $\psi(n+1)=H_n-\gamma$, which is the bridge between digamma and harmonic sums. Differentiating again gives the polygammas: $\psi_1(n+1)=\zeta(2)-\sum_{k=1}^{n}k^{-2}$, and in general $\psi_m(1)=(-1)^{m+1}m!\,\zeta(m+1)$.

\begin{worked}
Take $a=\tfrac12$ in the representation. Then $\psi(\tfrac12)=\psi(1)-2=-\gamma-2\ln2$ by the reflection value, so
\[\int_0^1\frac{1-x^{-1/2}}{1-x}\dd x=\psi\Bigl(\frac12\Bigr)+\gamma=-2\ln2\approx-1.386294.\]
A rational argument needs Gauss's theorem, which reflection plus the recurrence generates for small denominators. At $a=\tfrac13$,
\[\int_0^1\frac{1-x^{-2/3}}{1-x}\dd x=\psi\Bigl(\frac13\Bigr)+\gamma=-\frac{3}{2}\ln3-\frac{\pi}{2\sqrt3}\approx-2.554818,\]
which is confirmed to twelve digits by quadrature. Similarly $\psi(\tfrac14)+\gamma=-\tfrac{\pi}{2}-3\ln2$.
\end{worked}

\begin{worked}[The parameter derivative in action]
Where digamma really earns its keep is as the derivative of a Beta function. Differentiate
$\int_0^1x^{p-1}\dd x=\tfrac1p$ in $p$ to get $\int_0^1x^{p-1}\ln x\dd x=-\tfrac{1}{p^{2}}$, the single most used identity in the whole of parameter differentiation. Integrating that back in $p$ between two values gives the Frullani-type result
\[\int_0^1\frac{x^{a}-x^{b}}{\ln x}\dd x=\ln\frac{a+1}{b+1},\]
so $\int_0^1\frac{x^{3}-x}{\ln x}\dd x=\ln2\approx0.693147$, confirmed by quadrature.

One level up, differentiate the Beta function itself:
\[\pdv{}{p}\mathrm{B}(p,q)=\mathrm{B}(p,q)\bigl[\psi(p)-\psi(p+q)\bigr],\]
which evaluates $\int_0^1x^{p-1}(1-x)^{q-1}\ln x\dd x$ in one line. At $p=q=1$ that reads $\int_0^1\ln x\dd x=\psi(1)-\psi(2)=-1$, the correct answer, which is the check to run whenever you write the formula down.
\end{worked}

\begin{trap}
Both halves of $\dfrac{1-x^{a-1}}{1-x}$ blow up at $x=1$ individually; only the combination converges, so you may not split the integral and evaluate the pieces. And $\psi$ at a general rational argument is not reachable from the recurrence alone: it needs Gauss's digamma theorem, which supplies the $\cot$ and $\ln$ terms visible in the $a=\tfrac13$ answer above. Quoting $\psi(\tfrac13)=-\gamma-3\ln3$ or any other guess without the $\pi\cot$ piece is a standard way to lose a problem.
\end{trap}

\begin{seealsob}The Beta function; Euler-Mascheroni $\gamma$; Zeta values and Apery's constant; Feynman differentiation.\end{seealsob}

\section{The polylogarithm ladder}

Digamma handles one logarithm against a power. Two or more logarithms, or a logarithm against $\tfrac{1}{1-x}$, produce something Gamma cannot express, and the right object is the polylogarithm. It is worth understanding structurally rather than as a list: $\Li_s$ is defined so that dividing by $x$ and integrating raises $s$ by one, which makes the whole family a ladder that you climb one rung per integration.

\tech{Polylogarithm ladder}{hard}
\cue $\dfrac{\ln(1-x)}{x}$, $\dfrac{\ln^{k}x}{1\pm x}$, $\dfrac{\Li_n(x)}{x}$, or a series $\sum x^{n}/n^{s}$ emerging from a termwise integration. Any integral on $[0,1]$ whose integrand is a logarithm over a linear denominator lands here.
\method The definition and the ladder step:
\[\Li_s(x)=\sum_{n\geq1}\frac{x^{n}}{n^{s}},\qquad \dvop{x}\Li_{s+1}(x)=\frac{\Li_s(x)}{x},\qquad \Li_1(x)=-\ln(1-x).\]
So $\int_0^1\frac{\Li_s(x)}{x}\dd x=\Li_{s+1}(1)=\zeta(s+1)$. At the endpoints, $\Li_s(1)=\zeta(s)$ and $\Li_s(-1)=-\eta(s)=-(1-2^{1-s})\zeta(s)$. The two reduction identities worth carrying are Landen,
\[\Li_2(x)+\Li_2(1-x)=\frac{\pi^{2}}{6}-\ln x\ln(1-x),\]
and inversion, $\Li_2(-x)+\Li_2(-\tfrac1x)=-\tfrac{\pi^{2}}{6}-\tfrac12\ln^{2}x$. Between them they move any dilogarithm argument into $[0,\tfrac12]$, where the series converges fast.

\begin{worked}
The base case is the whole idea in miniature:
\[\int_0^1\frac{\ln(1-x)}{x}\dd x=-\int_0^1\frac{\Li_1(x)}{x}\dd x=-\Li_2(1)=-\frac{\pi^{2}}{6}\approx-1.644934.\]
Climbing two rungs at once,
\[\int_0^1\frac{\Li_3(x)-\Li_3(-x)}{x}\dd x=\Li_4(1)-\Li_4(-1)=\zeta(4)+\eta(4)=\frac{15}{8}\zeta(4)=\frac{\pi^{4}}{48}\approx2.029356,\]
since $\eta(4)=(1-2^{-3})\zeta(4)=\tfrac78\zeta(4)$.

A harder one, which is the standard route to $\zeta(3)$. Expand $\ln(1-x)=-\sum x^{n}/n$ in
\[\int_0^1\frac{\ln x\,\ln(1-x)}{x}\dd x=-\sum_{n\geq1}\frac1n\int_0^1x^{n-1}\ln x\dd x=\sum_{n\geq1}\frac{1}{n^{3}}=\zeta(3)\approx1.202057,\]
using $\int_0^1x^{n-1}\ln x\dd x=-1/n^{2}$. Each of the two logarithms contributed one power of $n$ downstairs; that is the arithmetic of the ladder.

Change the sign inside the logarithm and the same expansion produces the alternating counterparts:
\[\int_0^1\frac{\ln(1+x)}{x}\dd x=\eta(2)=\frac{\pi^{2}}{12}\approx0.822467,\qquad
\int_0^1\frac{\ln x\,\ln(1+x)}{x}\dd x=-\eta(3)=-\frac34\zeta(3)\approx-0.901543.\]
Both are verified numerically to twelve digits. The rule to take away: a $\ln(1-x)$ gives $\zeta$, a $\ln(1+x)$ gives $\eta$, and each additional $\ln x$ raises the order by one.
\end{worked}

\begin{trap}
Two sign traps. $\Li_s$ has a branch cut along $[1,\infty)$, and the inversion identity changes form across it, so applying the real version to $x>1$ produces a real answer for something with an imaginary part. And $\Li_n(-1)=-\eta(n)=-(1-2^{1-n})\zeta(n)$: the minus sign and the $\eta$-to-$\zeta$ conversion are separate opportunities to slip. Test the conversion at $n=2$, where it must give $-\pi^{2}/12$.
\end{trap}

\begin{seealsob}Zeta values and Apery's constant; Log-sine integrals and the Clausen function; Catalan's constant $G$; Geometric expansion.\end{seealsob}

\begin{keynote}[How to actually compute a dilogarithm]
The defining series $\sum x^{n}/n^{2}$ converges for $\abs x\leq1$ but slowly near $x=1$, where you need thousands of terms for six digits. The standard recipe is to move the argument first. If $x>\tfrac12$, use Landen, $\Li_2(x)=\tfrac{\pi^{2}}{6}-\ln x\ln(1-x)-\Li_2(1-x)$, which sends $x$ to $1-x$ and therefore into $[0,\tfrac12)$. If $\abs x>1$, use inversion first to get inside the unit disc. Then sum the series, which on $[0,\tfrac12]$ gains at least one digit every three terms. Check: $\Li_2(0.9)=1.2997147230$ by direct summation, and Landen reproduces it from $\Li_2(0.1)$, which needs eight terms.
\end{keynote}

The values at the top of the ladder are zeta values, and they appear so often that they deserve their own recognition cue. The pattern to internalise is that a $k$-fold logarithm against a geometric denominator produces $\zeta(k+1)$, with a rational or $\eta$-type factor in front decided by the sign in the denominator.

\tech[Zeta values and Aperys constant]{Zeta values and Apery's constant}{medium}
\cue $\dfrac{x^{s-1}}{\eu^{x}-1}$ or $\dfrac{x^{s-1}}{\eu^{x}+1}$ on $(0,\infty)$; $\dfrac{\ln^{k}x}{1\pm x}$ on $[0,1]$; a double or triple sum with denominators $n^{2}m$ or $n^{3}$; a Planck-type or Fermi-type integrand out of physics.
\method Expand the denominator as a geometric series and integrate termwise. The two master formulas are
\[\int_0^{\infty}\frac{x^{s-1}}{\eu^{x}-1}\dd x=\Gamma(s)\zeta(s),\qquad
\int_0^{\infty}\frac{x^{s-1}}{\eu^{x}+1}\dd x=\Gamma(s)\,\eta(s)=\bigl(1-2^{1-s}\bigr)\Gamma(s)\zeta(s),\]
both valid for $s>1$. The even values are elementary in $\pi$: $\zeta(2)=\tfrac{\pi^{2}}{6}$, $\zeta(4)=\tfrac{\pi^{4}}{90}$, $\zeta(6)=\tfrac{\pi^{6}}{945}$. The odd values are not. Apery's constant $\zeta(3)=1.2020569032\ldots$ is irrational and has no known expression in $\pi$ and elementary constants; when it appears, it is the answer.

\begin{worked}
The Planck integral, and its two logarithmic cousins:
\[\int_0^{\infty}\frac{x^{2}}{\eu^{x}-1}\dd x=\Gamma(3)\zeta(3)=2\zeta(3)\approx2.404114.\]
Substituting $x=-\ln t$ in the same integral turns it into $\int_0^1\frac{\ln^{2}t}{1-t}\dd t$, which therefore also equals $2\zeta(3)$. Flipping the sign in the denominator brings in $\eta$:
\[\int_0^1\frac{\ln^{2}t}{1+t}\dd t=\Gamma(3)\eta(3)=2\Bigl(1-\frac14\Bigr)\zeta(3)=\frac32\zeta(3)\approx1.803085,\]
and splitting into even and odd $n$ instead gives $\int_0^1\frac{\ln^{2}t}{1-t^{2}}\dd t=\tfrac74\zeta(3)\approx2.103600$. All three are the same expansion with different bookkeeping.
\end{worked}

\begin{trap}
The factor $\Gamma(s)$ is easy to drop, because at $s=2$ it equals $1$ and the formula looks like a pure zeta value. At $s=3$ it is $2$ and at $s=4$ it is $6$, and dropping it halves or sixths the answer. Separately: $\zeta(s)$ diverges at $s=1$, so the representation fails there; $\int_0^{\infty}\frac{\dd x}{\eu^{x}-1}$ is divergent at the origin even though the $\eta$ version, $\int_0^{\infty}\frac{\dd x}{\eu^{x}+1}=\ln2$, is perfectly finite.
\end{trap}

\begin{seealsob}Polylogarithm ladder; Gamma-function moment; Digamma and polygamma; Geometric expansion.\end{seealsob}

\begin{keynote}[Transcendental weight]
There is a grading on these constants that predicts the shape of an answer before you compute it. Assign weight $1$ to $\pi$ and to $\ln2$, weight $2$ to $\pi^{2}$, $\ln^{2}2$ and $\Li_2$ values, weight $3$ to $\zeta(3)$, $\pi^{2}\ln2$ and $\ln^{3}2$, and in general weight $k$ to $\zeta(k)$ and to any product whose weights sum to $k$. An integral on $[0,1]$ whose integrand carries $k$ logarithms almost always evaluates to a rational combination of weight-$k$ constants and nothing else. So $\int_0^1\frac{\ln^{2}x}{1+x}\dd x$ must be a multiple of $\zeta(3)$, and it is; $\int_0^{\pi/2}\ln^{2}\sin x\dd x$ must be a combination of $\pi\ln^{2}2$ and $\pi^{3}$, and it is. Use the grading as a filter on candidate answers: an answer of mixed weight is nearly always an algebra error.
\end{keynote}

\section{The Gaussian family, and what happens on a finite interval}

Gaussian integrals over the whole line are exact; over a finite interval they are not, and the closed form is the error function. Understanding that pair correctly is what stops you wasting time. The same phenomenon repeats one level up with oscillatory Gaussians, where the answer over the whole half-line is exact and the Fresnel integrals are the finite-interval version.

\tech{Gaussian and error function}{core}
\cue $\eu^{-x^{2}}$, or any $\eu^{-ax^{2}+bx+c}$, on the whole line, the half-line, or a finite interval. Also $\eu^{-x^{2}/2}$ with a probabilistic dressing, which is the same integral with $a=\tfrac12$.
\method Over the whole line the value is exact:
\[\int_{-\infty}^{\infty}\eu^{-ax^{2}}\dd x=\sqrt{\frac{\pi}{a}},\qquad \int_0^{\infty}\eu^{-x^{2}}\dd x=\frac{\sqrt\pi}{2},\]
the first proved by squaring and passing to polar coordinates. A linear term is removed by completing the square, which costs a constant factor and a shift of the (infinite) interval:
\[\int_{-\infty}^{\infty}\eu^{-ax^{2}+bx}\dd x=\sqrt{\frac{\pi}{a}}\;\eu^{b^{2}/(4a)}.\]
Over a finite interval there is no elementary antiderivative and the closed form is $\erf(z)=\tfrac{2}{\sqrt\pi}\int_0^{z}\eu^{-t^{2}}\dd t$, with $\erf(\infty)=1$ and $\erfc=1-\erf$.

\begin{worked}
Complete the square with $a=2$, $b=3$:
\[\int_{-\infty}^{\infty}\eu^{-2x^{2}+3x}\dd x=\sqrt{\frac{\pi}{2}}\,\eu^{9/8}\approx3.860479,\]
confirmed by quadrature to fifteen digits. On a finite interval the answer stops being elementary and starts being $\erf$:
\[\int_0^1\eu^{-x^{2}}\dd x=\frac{\sqrt\pi}{2}\erf(1)\approx0.746824,\]
and there is nothing further to do; $\erf(1)=0.842700793$ is a tabulated number. A power weight sends you back to $\Gamma$ rather than $\erf$: $\int_0^{\infty}\sqrt x\,\eu^{-x^{2}}\dd x=\tfrac12\Gamma(\tfrac34)\approx0.612708$.
\end{worked}

\begin{worked}[A Gaussian against an oscillation]
Multiply the Gaussian by a cosine and the answer is still a Gaussian, in the frequency. Set $F(b)=\int_0^{\infty}\eu^{-x^{2}}\cos(bx)\dd x$ and differentiate under the integral: $F'(b)=-\int_0^{\infty}x\eu^{-x^{2}}\sin(bx)\dd x$, which integrates by parts (with $u=\sin bx$, $\dd v=x\eu^{-x^{2}}\dd x$) to $-\tfrac b2F(b)$. That is a first-order linear ODE with $F(0)=\tfrac{\sqrt\pi}{2}$, so
\[\int_0^{\infty}\eu^{-x^{2}}\cos(bx)\dd x=\frac{\sqrt\pi}{2}\eu^{-b^{2}/4},\]
which at $b=1.3$ gives $0.5808386697$, matching quadrature to fifteen digits. The Fourier transform of a Gaussian is a Gaussian, and this is the shortest honest proof of it.
\end{worked}

\begin{trap}
Hunting for an elementary antiderivative on a finite interval. There is none, by Liouville's theorem, and $\erf$ \emph{is} the closed form: writing it is finishing, not giving up. The other standard slip is the factor of two: $\int_{-\infty}^{\infty}=\sqrt{\pi/a}$ while $\int_0^{\infty}$ is half that. And when completing the square, the shift $x\mapsto x-\tfrac{b}{2a}$ is harmless only because the interval is infinite; on a finite interval it moves the endpoints and the answer becomes a difference of two $\erf$ values.
\end{trap}

\begin{seealsob}Gamma-function moment; Fresnel integrals; The parameter ODE; Rotating the contour without contours.\end{seealsob}

Replace the decaying exponential $\eu^{-x^{2}}$ by the oscillating $\eu^{\iu x^{2}}$ and the integral still converges over the half-line, but for a completely different reason: not decay but cancellation. The values are the Fresnel integrals, and they are the first place in this chapter where convergence is only conditional, which changes what manipulations are legal.

\tech{Fresnel integrals}{hard}
\cue $\sin(x^{2})$ or $\cos(x^{2})$ on $(0,\infty)$; more generally $\sin(x^{p})$ for $p>1$; or $\dfrac{\sin u}{\sqrt u}$, which is what $u=x^{2}$ produces.
\method Substitute $u=x^{p}$ to reduce to a Mellin transform of the sine, then quote
\[\int_0^{\infty}u^{s-1}\sin u\dd u=\Gamma(s)\sin\frac{\pi s}{2},\qquad
\int_0^{\infty}u^{s-1}\cos u\dd u=\Gamma(s)\cos\frac{\pi s}{2},\qquad 0<s<1.\]
For $p=2$ this gives the classical pair
\[\int_0^{\infty}\sin(x^{2})\dd x=\int_0^{\infty}\cos(x^{2})\dd x=\frac12\sqrt{\frac{\pi}{2}}\approx0.626657.\]
The general statement is $\int_0^{\infty}\sin(x^{p})\dd x=\Gamma\bigl(1+\tfrac1p\bigr)\sin\tfrac{\pi}{2p}$.

\begin{worked}
Take $p=2$ and do it by hand. With $u=x^{2}$, $\dd x=\tfrac{\dd u}{2\sqrt u}$, so
\[\int_0^{\infty}\sin(x^{2})\dd x=\frac12\int_0^{\infty}u^{-1/2}\sin u\dd u=\frac12\Gamma\Bigl(\frac12\Bigr)\sin\frac{\pi}{4}=\frac{\sqrt\pi}{2}\cdot\frac{1}{\sqrt2}=\frac12\sqrt{\frac{\pi}{2}},\]
matching $0.6266571$ to seven digits under accelerated numerical summation of the half-period pieces.

For $p=3$ the same formula gives $\int_0^{\infty}\sin(x^{3})\dd x=\Gamma(\tfrac43)\sin\tfrac{\pi}{6}=\tfrac12\Gamma(\tfrac43)\approx0.446490$, again confirmed numerically.
\end{worked}

\begin{trap}
These converge only conditionally, so every step must respect that. No absolute-value comparison test, no naive Tonelli, no interchanging a limit with the integral without a separate argument. Practically: naive numerical quadrature on the raw $\sin(x^{2})$ over a truncated range does \emph{not} converge to the answer, because the tail contributes at the same order as your truncation error. Substitute $u=x^{2}$ first, or sum the alternating half-period contributions with an averaging acceleration, before you trust a numerical check.
\end{trap}

\begin{seealsob}Gaussian and error function; Rotating the contour without contours; Oscillatory convergence; Named-function answers.\end{seealsob}

\section{Oscillation divided by $x$: the sine-integral family}

The Fresnel integrals converge because the phase accelerates. The other standard mechanism is a slowly decaying amplitude against a fixed frequency, and its type specimen is $\frac{\sin x}{x}$. This is the most famous conditionally convergent integral in analysis, and its finite-interval version is a named function that many readers refuse to accept as an answer.

\tech[Dirichlet integral Si and Ci]{Dirichlet integral, $\Si$ and $\Ci$}{hard}
\cue $\dfrac{\sin x}{x}$ over an infinite range, or $\dfrac{\sin(\text{anything})}{\text{that same thing}}$ after a substitution; also bounds that are engineered to be absurdly large, which is a signal that the intended answer is a special-function value rather than a limit.
\method The Dirichlet value is
\[\int_0^{\infty}\frac{\sin x}{x}\dd x=\frac{\pi}{2},\]
derived by inserting a convergence factor: set $I(a)=\int_0^{\infty}\eu^{-ax}\frac{\sin x}{x}\dd x$, differentiate under the integral to get $I'(a)=-\frac{1}{1+a^{2}}$, integrate to $I(a)=\arctan\frac1a$, and let $a\to0^{+}$. For finite bounds, name the answer:
\[\Si(z)=\int_0^{z}\frac{\sin t}{t}\dd t,\qquad \Ci(z)=-\int_z^{\infty}\frac{\cos t}{t}\dd t=\gamma+\ln z+\int_0^{z}\frac{\cos t-1}{t}\dd t.\]

\begin{worked}
The substitution that manufactures a sine integral out of nothing. In $\int_{-\infty}^{20}\sin(\pi\eu^{x})\dd x$ put $u=\eu^{x}$, so $\dd x=\dd u/u$ and the range becomes $(0,\eu^{20})$:
\[\int_{-\infty}^{20}\sin(\pi\eu^{x})\dd x=\int_0^{\eu^{20}}\frac{\sin(\pi u)}{u}\dd u=\Si(\pi\eu^{20})\approx1.5707963270.\]
The special function is the closed form. The enormous upper limit is a joke at the reader's expense: $\Si(\pi\eu^{20})$ agrees with $\tfrac{\pi}{2}=1.5707963268$ to nine decimal places, but it is not equal to it.

The damping trick also handles the finite-$a$ family directly: $\int_0^{\infty}\eu^{-0.7x}\frac{\sin x}{x}\dd x=\arctan\frac{1}{0.7}\approx0.960070$.

Three relatives worth having by sight, all obtained by product-to-sum followed by the Dirichlet value:
\[\int_0^{\infty}\frac{\sin^{2}x}{x^{2}}\dd x=\frac{\pi}{2},\qquad
\int_0^{\infty}\frac{\sin^{3}x}{x}\dd x=\frac{\pi}{4},\qquad
\int_0^{\infty}\frac{\cos ax-\cos bx}{x}\dd x=\ln\frac ba.\]
The last is Frullani, and it is the cosine integral in disguise: the individual pieces are $\Ci$-divergent at the origin and only the difference survives, which is exactly the $\gamma$ mechanism of a later section. At $a=1$, $b=2$ it gives $\ln2\approx0.693147$.
\end{worked}

\begin{trap}
$\int_0^{\infty}\bigl|\frac{\sin x}{x}\bigr|\dd x=\infty$, so the integral is conditionally and not absolutely convergent, and dominated convergence is unavailable for the $a\to0^{+}$ step. That limit needs Abel's theorem or an explicit estimate on the oscillation, not a wave of the hand. A separate trap: $\Ci$ has a logarithmic singularity at the origin and $\Ci(z)\to-\infty$ as $z\to0^{+}$, so $\int_0^{z}\frac{\cos t}{t}\dd t$ is divergent while $\int_0^{z}\frac{\sin t}{t}\dd t$ is not. The two look symmetric and are not.
\end{trap}

\begin{seealsob}Lobachevsky's formula; Feynman differentiation; Fresnel integrals; Frullani's integral.\end{seealsob}

\begin{keynote}[Checking a conditionally convergent value]
The integrals of the last two techniques share a hazard: naive numerical quadrature on a truncated range does not converge to their values, because the tail contributes at the same order as the truncation. Two reliable fixes. Substitute first, so that $\int_0^{\infty}\sin(x^{2})\dd x$ becomes $\tfrac12\int_0^{\infty}u^{-1/2}\sin u\dd u$, where the amplitude decays. Or integrate over consecutive half-periods, forming an alternating series of the pieces, and accelerate it by repeated averaging; twenty averaging passes over two hundred half-periods gives twelve digits for $\int_0^{\infty}\frac{\sin x}{x}\dd x$ and for the Fresnel values. Every conditionally convergent number quoted in this chapter was checked that way, not by truncation.
\end{keynote}

\section{Logarithms of trigonometric functions}

The next family has nothing to do with oscillation. It arises whenever a logarithm meets a sine or cosine, which happens constantly, because $\ln\sin$ is what falls out of $\ln\abs{1-\eu^{\iu\theta}}$ and therefore out of any Fourier computation. The base value is elementary to derive and worth being able to reproduce in three lines.

\tech{Log-sine integrals and the Clausen function}{hard}
\cue $\ln\sin x$, $\ln\cos x$, $\ln\tan x$, $\ln\bigl(2\sin\tfrac x2\bigr)$, or a squared or cubed version, over $[0,\tfrac\pi2]$, $[0,\pi]$ or a partial range.
\method The base value comes from a fold. Writing $L=\int_0^{\pi/2}\ln\sin x\dd x$, complementarity gives $L=\int_0^{\pi/2}\ln\cos x\dd x$, so $2L=\int_0^{\pi/2}\ln(\tfrac12\sin2x)\dd x=-\tfrac{\pi}{2}\ln2+L$, hence
\[\int_0^{\pi/2}\ln\sin x\dd x=\int_0^{\pi/2}\ln\cos x\dd x=-\frac{\pi}{2}\ln2,\qquad \int_0^{\pi}\ln\sin x\dd x=-\pi\ln2.\]
The second moment needs the Fourier series $\ln(2\sin\tfrac\theta2)=-\sum_{n\geq1}\frac{\cos n\theta}{n}$ and gives
\[\int_0^{\pi/2}\ln^{2}(\sin x)\dd x=\frac{\pi}{2}\Bigl(\ln^{2}2+\frac{\pi^{2}}{12}\Bigr)\approx2.046622.\]
Partial ranges are not elementary: they define the Clausen function $\operatorname{Cl}_2(\theta)=-\int_0^{\theta}\ln\abs{2\sin\tfrac t2}\dd t=\sum_{n\geq1}\frac{\sin n\theta}{n^{2}}$.

\begin{worked}
The normalising $2$ is exactly what makes the full period vanish:
\[\int_0^{2\pi}\ln\Bigl(2\sin\frac x2\Bigr)\dd x=0,\]
which is the $n$-integrated form of the Fourier series above, every $\int_0^{2\pi}\cos n\theta\dd\theta$ being zero. Numerically the quadrature returns $-7\times10^{-15}$.

Two Clausen values worth knowing: $\operatorname{Cl}_2(\tfrac\pi2)=G=0.9159655942$, Catalan's constant, and the maximum $\operatorname{Cl}_2(\tfrac\pi3)=1.0149416064$, which is the largest value the function attains. And a mixed moment that combines this section with the last: $\int_0^{\pi/2}x\ln(\sin x)\dd x=\tfrac{7}{16}\zeta(3)-\tfrac{\pi^{2}}{8}\ln2\approx-0.329236$.

The cross moment is the cleanest illustration of the weight grading. Since $\ln\sin x\ln\cos x$ carries two logarithms, the answer must be a combination of $\pi\ln^{2}2$ and $\pi^{3}$, and it is:
\[\int_0^{\pi/2}\ln(\sin x)\ln(\cos x)\dd x=\frac{\pi}{2}\ln^{2}2-\frac{\pi^{3}}{48}\approx0.108730.\]
Expanding $\ln^{2}(\sin x\cos x)=\ln^{2}\sin x+2\ln\sin x\ln\cos x+\ln^{2}\cos x$ and using $\sin x\cos x=\tfrac12\sin2x$ reduces this to the two moments already computed.
\end{worked}

\begin{trap}
The integrand is $-\infty$ at the endpoint, so every manipulation is on an improper integral and the fold above needs the convergence checked once, at the start. The ``$2$'' inside $\ln(2\sin\tfrac x2)$ is not decorative: it is the normalisation that makes the mean value zero, and dropping it shifts every Clausen value by a multiple of $\ln2$. Finally, $\int_0^{\pi}\ln\sin x\dd x=-\pi\ln2$ is twice the quarter-range value, not the same as it; the two ranges give different numbers because the interval doubled, not because the integrand changed.
\end{trap}

\begin{seealsob}Catalan's constant $G$; Polylogarithm ladder; Fourier series kernels; Zeta values and Apery's constant.\end{seealsob}

\section{The constants that keep appearing}

There is a reason these three sit together rather than being scattered through the chapter. Each is a \emph{terminal} constant: an answer, not a step. Everything else in this chapter reduces, given enough identities, to $\pi$, $\ln2$, Gamma values at rational points, and these three. Knowing which numbers are terminal is what tells you when to stop, and stopping at the right moment is a large fraction of the skill.

Three numbers turn up again and again in this part of the subject, and each has a characteristic generating mechanism. Learning to recognise the mechanism is what lets you stop early instead of grinding for a closed form that does not exist. Catalan's constant comes from odd squares with alternating signs, Euler's $\gamma$ comes from the difference of two divergences, and $\zeta(3)$ comes from a threefold logarithmic weight. None of the three simplifies further.

\tech[Catalans constant]{Catalan's constant $G$}{medium}
\cue $\dfrac{\arctan x}{x}$ on $[0,1]$, $\dfrac{\ln x}{1+x^{2}}$ on $[0,1]$, $\int_0^{\pi/4}\ln(\cot x)\dd x$, $\dfrac{x}{\sin x}$ on $[0,\tfrac\pi2]$, or any sum that reduces to $\sum(-1)^{k}(2k+1)^{-2}$.
\method The constant is $G=\beta(2)=\sum_{k\geq0}\dfrac{(-1)^{k}}{(2k+1)^{2}}\approx0.9159655942$. It appears whenever an alternating odd-denominator series, that is an $\arctan$ or a $\Li$ at $\iu$, is integrated once more against $\tfrac{\dd x}{x}$ or $\ln x\dd x$, because $\int_0^1x^{2k}\dd x=\tfrac{1}{2k+1}$ supplies the second odd denominator. Equivalently $G=\operatorname{Cl}_2(\tfrac{\pi}{2})$.

\begin{worked}
Expand and integrate termwise:
\[\int_0^1\frac{\arctan x}{x}\dd x=\int_0^1\sum_{k\geq0}\frac{(-1)^{k}x^{2k}}{2k+1}\dd x=\sum_{k\geq0}\frac{(-1)^{k}}{(2k+1)^{2}}=G\approx0.915966.\]
The same constant, three more ways, all verified numerically to fifteen digits:
\[\int_0^1\frac{-\ln x}{1+x^{2}}\dd x=G,\qquad \int_0^{\pi/4}\ln(\cot x)\dd x=G,\qquad \int_0^{\pi/2}\frac{x}{\sin x}\dd x=2G.\]
The second is the first after $x=\tan\theta$; the third comes from the same $\arctan$ series against a different weight.

A case where $G$ appears with a companion, which is the usual situation: $\int_0^1\frac{\ln(1+x^{2})}{1+x^{2}}\dd x=\frac{\pi}{2}\ln2-G\approx0.172827$. Substituting $x=\tan\theta$ turns it into $-2\int_0^{\pi/4}\ln\cos\theta\dd\theta$, so it sits at the junction of this section and the log-sine section, and the two constants that appear are exactly the two that the weight grading permits.
\end{worked}

\begin{trap}
$G$ has no known expression in $\pi$, $\ln2$ and zeta values, and it is not known whether it is irrational. Grinding to reduce it further wastes the problem. Recognise it, write it, and stop. The converse trap is naming it too early: $\int_0^1\frac{\arctan x}{1+x^{2}}\dd x$ looks like the same family but is elementary, equal to $\tfrac12\arctan^{2}(1)=\tfrac{\pi^{2}}{32}$, because the weight is the derivative of the argument.
\end{trap}

\begin{seealsob}Log-sine integrals and the Clausen function; Polylogarithm ladder; Zeta values and Apery's constant.\end{seealsob}

Catalan's constant is generated by a series. The Euler-Mascheroni constant is generated by a subtraction, and that difference in mechanism is the whole reason it behaves differently: it appears in problems where two things separately diverge.

\tech[Euler-Mascheroni gamma]{Euler-Mascheroni $\gamma$}{medium}
\cue $\eu^{-x}\ln x$ on the half-line; $\ln\ln\tfrac1x$ on $[0,1]$; a bracketed pair of divergent terms such as $\dfrac{1}{\ln x}+\dfrac{1}{1-x}$; or the difference between a harmonic-type divergence and a logarithmic one.
\method The definition that matters for integration is $\gamma=-\Gamma'(1)=-\psi(1)$, that is,
\[\gamma=-\int_0^{\infty}\eu^{-x}\ln x\dd x\approx0.5772156649.\]
Two equivalent forms cover most sightings:
\[\int_0^1\ln\ln\frac1x\dd x=-\gamma,\qquad \int_0^1\Bigl(\frac{1}{\ln x}+\frac{1}{1-x}\Bigr)\dd x=\gamma,\]
the first being the previous one after $x=\eu^{-t}$, the second a genuine cancellation of two logarithmic divergences at $x=1$.

\begin{worked}
Verify the pair form directly. Near $x=1$ write $x=1-\varepsilon$: then $\tfrac{1}{\ln x}\approx-\tfrac1\varepsilon-\tfrac12$ and $\tfrac{1}{1-x}=\tfrac1\varepsilon$, so the $\tfrac1\varepsilon$ poles cancel and the integrand tends to $-\tfrac12$. The integral therefore converges, and quadrature gives $0.5772156649$, matching $\gamma$ to ten digits.

The second moment brings in $\zeta(2)$ alongside $\gamma$, and is worth knowing because it is how $\gamma^{2}$ enters problems at all:
\[\int_0^{\infty}\eu^{-x}\ln^{2}x\dd x=\gamma^{2}+\frac{\pi^{2}}{6}\approx1.978112,\]
which is $\Gamma''(1)$, and follows from $\Gamma''(1)=\Gamma(1)\bigl[\psi(1)^{2}+\psi_1(1)\bigr]$ with $\psi(1)=-\gamma$ and $\psi_1(1)=\zeta(2)$.

A cue in disguise: $\int_0^1\ln\ln\tfrac1x\dd x$. Put $x=\eu^{-t}$, so $\dd x=-\eu^{-t}\dd t$ and the integral becomes $\int_0^{\infty}\eu^{-t}\ln t\dd t=-\gamma\approx-0.577216$. Whenever $\ln\ln$ appears on $[0,1]$, that substitution is the move.
\end{worked}

\begin{trap}
Splitting the pair. Each of $\int_0^1\frac{\dd x}{\ln x}$ and $\int_0^1\frac{\dd x}{1-x}$ diverges logarithmically at $x=1$, so the difference is meaningful only if the two cut-offs are taken \emph{simultaneously}: regularise as $\int_0^{1-\varepsilon}$ of the whole bracket and then let $\varepsilon\to0$. Cutting the two terms at different points changes the answer by any constant you like.
\end{trap}

\begin{seealsob}Digamma and polygamma; Frullani's integral; Epsilon-regularisation of a divergent pair.\end{seealsob}

\section{Where the ladder runs out}

Everything so far has been, in the end, Gamma and its derivatives plus the polylogarithm. That covers a startling amount of ground, but it stops at a definite frontier: a square root of a cubic or quartic. There the integral defines a genuinely new function of the coefficients, and no rearrangement brings it back.

\tech[Elliptic integrals K and E]{Elliptic integrals $K$ and $E$}{hard}
\cue $\dfrac{1}{\sqrt{1-k^{2}\sin^{2}\theta}}$; $\dfrac{1}{\sqrt{(1-x^{2})(1-k^{2}x^{2})}}$; $\dfrac{1}{\sqrt{1-x^{4}}}$ or any square root of a cubic or quartic with distinct roots; the exact period of a pendulum at finite amplitude; the arc length of an ellipse.
\method Name them and stop:
\[K(k)=\int_0^{\pi/2}\frac{\dd\theta}{\sqrt{1-k^{2}\sin^{2}\theta}},\qquad E(k)=\int_0^{\pi/2}\sqrt{1-k^{2}\sin^{2}\theta}\dd\theta.\]
A quartic under a square root is elliptic unless a special substitution reduces it. At the lemniscatic value $k^{2}=\tfrac12$ the elliptic integral does collapse to Gammas, which is why $\int_0^1(1-x^{4})^{-1/2}\dd x$ has a closed form while the general case does not.

\begin{worked}
The lemniscatic case, done through Beta rather than through elliptic tables. Put $u=x^{4}$:
\[\int_0^1\frac{\dd x}{\sqrt{1-x^{4}}}=\frac14\int_0^1u^{-3/4}(1-u)^{-1/2}\dd u=\frac14\mathrm{B}\Bigl(\frac14,\frac12\Bigr)=\frac{\Gamma(\tfrac14)\Gamma(\tfrac12)}{4\Gamma(\tfrac34)}.\]
Now use reflection, $\Gamma(\tfrac34)=\pi\sqrt2/\Gamma(\tfrac14)$:
\[=\frac{\Gamma(\tfrac14)^{2}\sqrt\pi}{4\pi\sqrt2}=\frac{\Gamma(\tfrac14)^{2}}{4\sqrt{2\pi}}\approx1.311029,\]
matching quadrature to twelve digits. Equivalently this is $\tfrac{1}{\sqrt2}K(\tfrac{1}{\sqrt2})$ with $K(\tfrac{1}{\sqrt2})=\Gamma(\tfrac14)^{2}/(4\sqrt\pi)\approx1.854075$.
\end{worked}

\begin{worked}[The pendulum, which is where these came from]
A pendulum of length $L$ released from angle $\theta_0$ has exact period
\[T=4\sqrt{\frac Lg}\,K\Bigl(\sin\frac{\theta_0}{2}\Bigr),\]
against the small-angle prediction $2\pi\sqrt{L/g}$. At $\theta_0=\tfrac\pi2$ the modulus is $\sin\tfrac\pi4=\tfrac{1}{\sqrt2}$, so the ratio of true period to small-angle period is
\[\frac{4K(1/\sqrt2)}{2\pi}=\frac{2\cdot1.8540746773}{\pi}\approx1.1803406,\]
an 18 percent correction. The lemniscatic modulus is not a coincidence here: $\theta_0=\tfrac\pi2$ is exactly the amplitude at which the elliptic integral collapses to Gammas.
\end{worked}

\begin{trap}
Two conventions for the modulus, and software disagrees with tables. Some write $K(k)$ with $k$ the modulus, others $K(m)$ with $m=k^{2}$ the parameter; most numerical libraries use $m$. State which you use, and sanity-check against $K$ at a value you know. The other trap runs the other way: a cubic or quartic radical is not automatically elliptic. If the radical is $\sqrt{Q(x)}$ and the rest of the integrand is a multiple of $Q'(x)$, the integral is elementary, because $\int Q'\sqrt{Q}\dd x=\tfrac23Q^{3/2}$. Look for that before you name anything.
\end{trap}

\begin{seealsob}The Beta function; Euler reflection and the Mellin transform; Named-function answers; Sine substitution.\end{seealsob}

Elliptic integrals are the best-known place where naming is the answer, but they are not the only one. The general habit is worth stating explicitly, because it is the habit this whole chapter is trying to install: when a substitution lands you on a shape that has a name, use the name and then use its identities, rather than continuing to push symbols around.

\tech{Named-function answers}{hard}
\cue A substitution lands you on $w\eu^{w}$, on $\cos(\sin x)$, on $\int_0^{z}\eu^{-t}t^{a-1}\dd t$, or on an oscillatory kernel with a $\cosh$ or an Airy function in the denominator. In general: a shape that resists every elementary method but keeps reappearing.
\method Stop and name it. Lambert $W$ is defined by $W(x)\eu^{W(x)}=x$; the Bessel function by $J_0(z)=\tfrac1\pi\int_0^{\pi}\cos(z\sin x)\dd x$; the incomplete Gammas by $\gamma(a,z)=\int_0^{z}\eu^{-t}t^{a-1}\dd t$ and $\Gamma(a,z)=\Gamma(a)-\gamma(a,z)$; and the Airy pair through the Wronskian identity $\dvop{z}\arctan\frac{\operatorname{Bi}(z)}{\operatorname{Ai}(z)}=\frac{1/\pi}{\operatorname{Ai}^{2}+\operatorname{Bi}^{2}}$. Then integrate the named object using its own identities, most often inverse-function integration or a recurrence.

\begin{worked}
Two instances. First, recognition of a Bessel integral:
\[\int_0^{\pi}\cos(\sin x)\dd x=\pi J_0(1)\approx2.403939,\]
straight from the definition with $z=1$. Second, an inverse function integrated by the graph-swapping identity: with $x=t\eu^{t}$, so that $W(x)=t$ and $\dd x=(1+t)\eu^{t}\dd t$, the limits $x=0$ and $x=\eu$ become $t=0$ and $t=1$, and
\[\int_0^{\eu}W(x)\dd x=\int_0^1t(1+t)\eu^{t}\dd t=\eu-1\approx1.718282.\]
The remaining integral is elementary because the substitution moved the difficulty into the limits.
\end{worked}

\begin{worked}[Incomplete Gammas, and why they are not a cop-out]
Truncating a Gamma integral gives the incomplete Gamma, and it obeys the same recurrence as $\Gamma$ with a boundary term:
\[\gamma(a+1,z)=a\gamma(a,z)-z^{a}\eu^{-z}.\]
So $\int_0^1\sqrt x\,\eu^{-x}\dd x=\gamma(\tfrac32,1)=\tfrac12\gamma(\tfrac12,1)-\eu^{-1}$, and $\gamma(\tfrac12,1)=\sqrt\pi\erf(1)$, giving
\[\int_0^1\sqrt x\,\eu^{-x}\dd x=\frac{\sqrt\pi}{2}\erf(1)-\frac1\eu\approx0.3789447,\]
confirmed by quadrature. Notice that the incomplete Gamma reduced to $\erf$, which is the general pattern: at half-integer $a$ the incomplete Gammas are error functions, and at integer $a$ they are finite sums times $\eu^{-z}$.
\end{worked}

\begin{trap}
The obvious failure is continuing to hunt for an elementary form. The subtler one runs the other way: naming an object when a cancellation was available. Two Bessel terms of equal magnitude and opposite sign cancel exactly, and a difference of incomplete Gammas often telescopes to something elementary. Evaluate and simplify before you name, then name and stop.
\end{trap}

\begin{seealsob}Elliptic integrals $K$ and $E$; Inverse-function integration; Dirichlet integral, $\Si$ and $\Ci$; Disguised zero.\end{seealsob}

\section{Three problems worked end to end}

Individual techniques are easy to demonstrate and easy to overrate. What the chapter is really teaching is the chaining, so here are three complete problems in which no single step is difficult and the whole is well past what elementary methods reach. Read them for the order of the moves, not for the algebra.

The first shows parameter differentiation applied to a special-function formula rather than to an integral. Consider
\[I=\int_0^{\infty}\frac{\ln x}{1+x^{3}}\dd x.\]
There is no obvious substitution and no symmetry, but there is a family: the reflection master formula gives $\int_0^{\infty}\frac{x^{s-1}}{1+x^{n}}\dd x=\frac{\pi}{n\sin(\pi s/n)}$ for the whole strip $0<s<n$, and $\ln x$ is what $\partial_s$ of $x^{s-1}$ produces. Differentiate both sides in $s$:
\[\int_0^{\infty}\frac{x^{s-1}\ln x}{1+x^{n}}\dd x=-\frac{\pi^{2}}{n^{2}}\cdot\frac{\cos(\pi s/n)}{\sin^{2}(\pi s/n)}.\]
Set $s=1$, $n=3$. Then $\cos\tfrac\pi3=\tfrac12$ and $\sin^{2}\tfrac\pi3=\tfrac34$, so
\[I=-\frac{\pi^{2}}{9}\cdot\frac{1/2}{3/4}=-\frac{2\pi^{2}}{27}\approx-0.7310818075,\]
which quadrature confirms to twelve digits. The same line at $n=4$ gives $\int_0^{\infty}\frac{\ln x}{1+x^{4}}\dd x=-\frac{\pi^{2}\sqrt2}{16}\approx-0.872358$. Notice what made this work: not cleverness about the integrand, but the fact that the answer was already known as a function of a parameter, and $\ln x$ is a parameter derivative.

The second problem chains Beta, digamma and trigamma to reach a log-sine moment that was quoted without proof earlier. Start from the trigonometric Beta integral in its one-parameter form,
\[f(p)=\int_0^{\pi/2}\sin^{2p-1}x\dd x=\frac12\mathrm{B}\Bigl(p,\frac12\Bigr)=\frac{\sqrt\pi\,\Gamma(p)}{2\Gamma(p+\tfrac12)}.\]
Differentiating in $p$ brings down a factor $2\ln\sin x$ each time, so that
\[f'(p)=2\int_0^{\pi/2}\sin^{2p-1}x\,\ln\sin x\dd x,\qquad f''(p)=4\int_0^{\pi/2}\sin^{2p-1}x\,\ln^{2}\sin x\dd x.\] The right-hand side is easiest to differentiate logarithmically:
\[\frac{f'}{f}=\psi(p)-\psi\Bigl(p+\frac12\Bigr),\qquad \frac{f''}{f}=\Bigl[\psi(p)-\psi\Bigl(p+\frac12\Bigr)\Bigr]^{2}+\psi_1(p)-\psi_1\Bigl(p+\frac12\Bigr).\]
Evaluate at $p=\tfrac12$, where $f(\tfrac12)=\tfrac\pi2$, $\psi(\tfrac12)-\psi(1)=-2\ln2$, and $\psi_1(\tfrac12)-\psi_1(1)=\tfrac{\pi^{2}}{2}-\tfrac{\pi^{2}}{6}=\tfrac{\pi^{2}}{3}$. The first derivative gives
\[\int_0^{\pi/2}\ln\sin x\dd x=\frac{f'(\tfrac12)}{2}=\frac{1}{2}\cdot\frac{\pi}{2}(-2\ln2)=-\frac{\pi}{2}\ln2,\]
recovering the base value, and the second gives
\[\int_0^{\pi/2}\ln^{2}\sin x\dd x=\frac{f''(\tfrac12)}{4}=\frac14\cdot\frac{\pi}{2}\Bigl(4\ln^{2}2+\frac{\pi^{2}}{3}\Bigr)=\frac{\pi}{2}\ln^{2}2+\frac{\pi^{3}}{24}\approx2.046622.\]
That is exactly the value quoted earlier, now derived. Four named objects appeared, and each was used for precisely one thing: Beta to get a closed form with a parameter, digamma to differentiate it once, trigamma to differentiate it twice, and the reflection value $\psi(\tfrac12)=-\gamma-2\ln2$ to evaluate at the point of interest.

The third problem is the second derivative of the first, and it is included to make a point about how much mileage one parametrised formula gives. Differentiating the master formula twice in $s$,
\[\int_0^{\infty}\frac{x^{s-1}\ln^{2}x}{1+x^{n}}\dd x=\dv{^{2}}{s^{2}}\Bigl[\frac{\pi}{n}\csc\frac{\pi s}{n}\Bigr]
=\frac{\pi^{3}}{n^{3}}\Bigl[\csc^{3}\frac{\pi s}{n}+\csc\frac{\pi s}{n}\cot^{2}\frac{\pi s}{n}\Bigr].\]
At $s=1$, $n=2$ the cotangent vanishes and the cosecant is $1$, so
\[\int_0^{\infty}\frac{\ln^{2}x}{1+x^{2}}\dd x=\frac{\pi^{3}}{8}\approx3.8757845850,\]
confirmed to eleven digits. The first derivative at the same point vanishes, which recovers the familiar $\int_0^{\infty}\frac{\ln x}{1+x^{2}}\dd x=0$, a fact usually proved by the inversion $x\mapsto1/x$; here it falls out as a stationary point of $\csc$. One formula, one parameter, and three separate results, two of which look like different problems entirely.

\begin{keynote}[The four errors that account for most failures]
Collected here because they are worth rereading before an examination. One: applying an identity outside its strip of validity, which yields a plausible finite number for a divergent integral. Two: confusing the algebraic and trigonometric Beta conventions, $p-1$ against $2p-1$. Three: dropping the $\Gamma(s)$ factor in $\int_0^{\infty}\frac{x^{s-1}}{\eu^{x}-1}\dd x=\Gamma(s)\zeta(s)$, invisible at $s=2$ and wrong by a factor $2$ at $s=3$. Four: splitting a bracketed pair of divergences, or splitting $\frac{1-x^{a-1}}{1-x}$, where only the combination converges. All four produce clean, confident, wrong answers, and all four are caught in thirty seconds by a numerical check.
\end{keynote}

\section{A table of values, and how to choose a rung}

Nearly every problem in this family is finished by one of a short list of numbers. Keeping the list where you can see it converts a hard integral into a lookup, and, more importantly, lets you recognise a target: if a numerical estimate of your integral is $0.9159655942$, you know what the answer is and can work backwards to a proof.

\begin{center}
\footnotesize
\setlength{\tabcolsep}{5pt}
\renewcommand{\arraystretch}{1.3}
\begin{tabular}{@{}lll@{}}
\toprule
\mh{Quantity} & \mh{Exact value} & \mh{Decimal} \\
\midrule
$\Gamma(\tfrac12)$ & $\sqrt\pi$ & $1.7724538509$ \\
$\Gamma(\tfrac13)$ & no simpler form & $2.6789385347$ \\
$\Gamma(\tfrac14)$ & no simpler form & $3.6256099082$ \\
$\Gamma(\tfrac34)$ & $\pi\sqrt2/\Gamma(\tfrac14)$ & $1.2254167025$ \\
$\Gamma(\tfrac13)\Gamma(\tfrac23)$ & $2\pi/\sqrt3$ & $3.6275987285$ \\
$\psi(1)$ & $-\gamma$ & $-0.5772156649$ \\
$\psi(\tfrac12)$ & $-\gamma-2\ln2$ & $-1.9635100260$ \\
$\psi(\tfrac13)$ & $-\gamma-\tfrac32\ln3-\tfrac{\pi}{2\sqrt3}$ & $-3.1320337800$ \\
$\psi(\tfrac14)$ & $-\gamma-\tfrac{\pi}{2}-3\ln2$ & $-4.2274535334$ \\
$\zeta(2)$ & $\pi^{2}/6$ & $1.6449340668$ \\
$\zeta(3)$ & Apery's constant & $1.2020569032$ \\
$\zeta(4)$ & $\pi^{4}/90$ & $1.0823232337$ \\
$\eta(2)$ & $\pi^{2}/12$ & $0.8224670334$ \\
$\Li_2(\tfrac12)$ & $\tfrac{\pi^{2}}{12}-\tfrac12\ln^{2}2$ & $0.5822405265$ \\
$G=\operatorname{Cl}_2(\tfrac\pi2)$ & no simpler form & $0.9159655942$ \\
$\operatorname{Cl}_2(\tfrac\pi3)$ & maximum of $\operatorname{Cl}_2$ & $1.0149416064$ \\
$\gamma$ & $-\Gamma'(1)$ & $0.5772156649$ \\
$\erf(1)$ & $\tfrac{2}{\sqrt\pi}\int_0^1\eu^{-t^{2}}$ & $0.8427007929$ \\
$\Si(\infty)$ & $\pi/2$ & $1.5707963268$ \\
$K(\tfrac{1}{\sqrt2})$ & $\Gamma(\tfrac14)^{2}/(4\sqrt\pi)$ & $1.8540746773$ \\
$\int_0^{\infty}\sin(x^{2})\dd x$ & $\tfrac12\sqrt{\pi/2}$ & $0.6266570687$ \\
\bottomrule
\end{tabular}
\end{center}

The choice of rung is decided almost entirely by the integrand's ingredients, not by the interval, which is the reverse of the situation in the symmetry chapter. Powers and exponentials mean Gamma; powers on a bounded interval mean Beta; a logarithm means differentiate a parameter and expect digamma or a polylogarithm; oscillation means Fresnel or the sine integral; a quartic radical means elliptic.

\begin{center}
\begin{tikzpicture}[node distance=6mm and 7mm]
\node[dtest] (a) {What is the hardest\\ ingredient?};
\node[dtest, below left=9mm and 6mm of a] (b) {Power against\\ $\eu^{-x^{k}}$ or $(1-x)^{q}$?};
\node[dtest, below right=9mm and 2mm of a] (c) {Logarithm, or\\ oscillation?};
\node[dleaf, below left=9mm and -2mm of b] (d) {Half-line: $\Gamma$\\ Unit interval: $\mathrm{B}$\\ Fractional power: reflect};
\node[dleaf, below right=9mm and -4mm of b] (e) {Quartic radical:\\ $K$, $E$, or Beta at\\ the lemniscatic point};
\node[dleaf, below left=9mm and -4mm of c] (f) {One $\ln$: $\psi$\\ Several: $\Li_k$, $\zeta$\\ $\ln\sin$: Clausen};
\node[dleaf, below right=9mm and -2mm of c] (g) {$\sin(x^{2})$: Fresnel\\ $\sin x/x$: $\Si$\\ $\eu^{-x^{2}}$: $\erf$};
\draw[dedge] (a) -- node[dlab,left]{algebraic} (b);
\draw[dedge] (a) -- node[dlab,right]{transcendental} (c);
\draw[dedge] (b) -- node[dlab,left]{yes} (d);
\draw[dedge] (b) -- node[dlab,right]{no} (e);
\draw[dedge] (c) -- node[dlab,left]{$\ln$} (f);
\draw[dedge] (c) -- node[dlab,right]{oscillation} (g);
\end{tikzpicture}
\end{center}

\begin{keynote}[Guess the constant first]
Before doing any work, estimate the integral numerically to four or five digits and compare against the table. This is not cheating; it is the single most efficient move available at this level. If the estimate is $2.0293561$ you are looking at $\pi^{4}/48$, and knowing that tells you the answer is a fourth-order polylogarithm, which tells you to expand a logarithm and integrate termwise. If the estimate matches nothing, the answer is probably a product or ratio of table entries, and $\ln$ of the estimate is worth checking too. The characteristic mistake at this level is not a wrong technique but a correct technique aimed at the wrong target.
\end{keynote}

It is worth saying explicitly when this chapter is the wrong tool. A named function is the right answer only when the integral genuinely resists elementary evaluation. If the integrand is a derivative in disguise, if a symmetry fold closes it, or if a substitution makes it rational, none of the machinery here is needed and reaching for it will cost you time and introduce branch conditions you did not have to satisfy. The order of attack is always the same: simplify algebraically, test the standard folds, try the obvious substitutions, and only then ask which special function is hiding. The one exception is when a preliminary numerical estimate already matches a table entry, in which case you may as well aim straight at the target.

One closing warning about the whole chapter. The identities here are all conditional: reflection needs $0<s<1$, the Mellin transforms of $\sin$ and $\cos$ need $0<s<1$, the zeta representation needs $s>1$, and the polylogarithm identities need the argument off the branch cut. In elementary integration a violated hypothesis usually produces an obvious absurdity like a negative area. Here it produces a plausible number, and the only defence is to check the strip explicitly and to verify the final value numerically. Every closed form printed in this chapter was checked that way.

\section{Recognition matrix}
\begin{recog}{@{}p{0.31\textwidth}p{0.35\textwidth}p{0.28\textwidth}@{}}
\rowcolor{mist}\mh{If you see} & \mh{Reach for} & \mh{If that fails} \\
\midrule
\endhead
$x^{m}\eu^{-\lambda x^{k}}$ on $(0,\infty)$ & $u=\lambda x^{k}$; $\tfrac{1}{k\lambda^{(m+1)/k}}\Gamma(\tfrac{m+1}{k})$ & Check $\tfrac{m+1}{k}>0$ \\
$a^{-x^{2}}$ with $a\neq\eu$ & Write $\lambda=\ln a$, then Gamma & Watch the power of $\ln a$ \\
$\dfrac{x^{s-1}}{1+x^{n}}$ on $(0,\infty)$ & Reflection: $\dfrac{\pi}{n\sin(\pi s/n)}$ & Verify $0<s<n$ \\
$\Gamma(\alpha)\Gamma(1-\alpha)$ anywhere & $\pi/\sin\pi\alpha$ & Shift by the recurrence first \\
$\Gamma(z)$ and $\Gamma(z+\tfrac12)$ together & Duplication: $2^{1-2z}\sqrt\pi\,\Gamma(2z)$ & Test the constant at $z=\tfrac12$ \\
$x^{p-1}(1-x)^{q-1}$ on $[0,1]$ & $\mathrm{B}(p,q)=\Gamma(p)\Gamma(q)/\Gamma(p+q)$ & Substitute to reach $[0,1]$ \\
$\sin^{m}x\cos^{n}x$, non-integer & $\tfrac12\mathrm{B}(\tfrac{m+1}{2},\tfrac{n+1}{2})$ & Solve $2p-1=m$ explicitly \\
$x^{a}(1-x^{b})^{c}$ & $u=x^{b}$; $\tfrac1b\mathrm{B}(\tfrac{a+1}{b},c+1)$ & Then reflection for fractions \\
A single $\ln x$ against a power & Differentiate Beta or Gamma in the parameter & Expect $\psi$ and $\gamma$ \\
$\dfrac{1-x^{a-1}}{1-x}$ on $[0,1]$ & $\psi(a)+\gamma$ & Gauss's digamma theorem \\
Harmonic numbers in a series & $\psi(n+1)=H_n-\gamma$ & Abel summation \\
$\dfrac{\ln(1-x)}{x}$ or $\dfrac{\Li_n(x)}{x}$ & Climb the ladder: $\Li_{n+1}$ & Landen or inversion \\
$\dfrac{\ln^{k}x}{1\pm x}$ on $[0,1]$ & Expand; $k!\,\zeta(k+1)$ or $k!\,\eta(k+1)$ & Keep the $\Gamma(s)$ factor \\
$\dfrac{x^{s-1}}{\eu^{x}\mp1}$ on $(0,\infty)$ & $\Gamma(s)\zeta(s)$ or $\Gamma(s)\eta(s)$ & Diverges at $s=1$ for the minus \\
$\eu^{-ax^{2}+bx}$ on $\R$ & Complete the square; $\sqrt{\pi/a}\,\eu^{b^{2}/4a}$ & Halve it for $(0,\infty)$ \\
$\eu^{-x^{2}}$ on a finite interval & $\tfrac{\sqrt\pi}{2}\erf(z)$; that is the answer & No elementary form exists \\
$\sin(x^{p})$ or $\cos(x^{p})$ on $(0,\infty)$ & $u=x^{p}$; $\Gamma(1+\tfrac1p)\sin\tfrac{\pi}{2p}$ & Only conditionally convergent \\
$\dfrac{\sin x}{x}$ to $\infty$ & $\tfrac\pi2$; damp with $\eu^{-ax}$ to prove it & Finite bound: $\Si(z)$ \\
An absurdly large upper limit & The answer is a special-function value & Substitute to expose $\tfrac{\sin u}{u}$ \\
$\ln\sin x$ or $\ln\cos x$ on $[0,\tfrac\pi2]$ & Fold with $\sin2x$; $-\tfrac\pi2\ln2$ & Fourier series of $\ln(2\sin\tfrac\theta2)$ \\
$\ln\bigl(2\sin\tfrac x2\bigr)$ on a partial range & Clausen $\operatorname{Cl}_2(\theta)$ & Keep the factor $2$ \\
$\dfrac{\arctan x}{x}$ or $\dfrac{\ln x}{1+x^{2}}$ on $[0,1]$ & Catalan's constant $G$ & Stop; $G$ does not reduce \\
$\ln\ln\tfrac1x$ on $[0,1]$ & $x=\eu^{-t}$; the answer is $-\gamma$ & Compare with $-\Gamma'(1)$ \\
A bracketed pair of divergences & Regularise together; expect $\gamma$ & Never split the bracket \\
$\sqrt{\text{cubic}}$ or $\sqrt{\text{quartic}}$ & $K$ and $E$; name and stop & Check for $\int Q'\sqrt Q$ first \\
$w\eu^{w}$, $\cos(\sin x)$, $\operatorname{Ai}$ & Lambert $W$, $J_0$, Airy; use their identities & Simplify before naming \\
A numerical value you recognise & Look it up in the table, then work backwards & Try ratios and logs of it \\
\end{recog}
